%------------------------------------------------- Chapter 2: state of the art
\chapterseparator{2}{Market Study and State of the Art}

\chapter{Market Study and State of the Art}

\section{Introduction}

This chapter situates the project within its commercial, algorithmic, and methodological landscape. It begins with an analytical study of the e-commerce personalisation market, emphasizing the domain-specific constraints of the cosmetic and wellness sector. It then examines the fundamental concepts underpinning modern conversational systems: Natural Language Processing, Large Language Models, Retrieval-Augmented Generation (RAG), vector representations, tool calling, and Docker containerisation. It reviews existing academic and industrial solutions, establishes their limitations, formalizes the proposed hybrid architecture, and details the CRISP-DM methodology adopted throughout the project lifecycle.

\section{Market analysis}

\subsection{Personalisation in online retail}

Product recommendation has become an essential capability in modern digital commerce. When an online catalogue expands beyond what a visitor can easily browse, discovery becomes a friction point that depresses conversion rates. Personalisation shortens the journey from latent customer need to relevant product discovery, lowering cart abandonment while enhancing user engagement.

In e-commerce literature and industry practice, three distinct recommendation paradigms are recognized:
\begin{itemize}
  \item \textbf{Behavioural recommendation}: Based on collaborative purchase and navigation histories (``Customers who bought this item also bought...'');
  \item \textbf{Similarity-based recommendation}: Suggesting adjacent or alternative items with comparable specifications to the product currently viewed;
  \item \textbf{Advisory recommendation}: Guiding the customer through a consultative dialogue to translate symptom-level needs into an ordered product sequence.
\end{itemize}

While behavioural and item-similarity engines are widely commoditized, consultative advisory recommendation remains rare in small and medium enterprises due to high integration costs and data sparsity.

\subsection{Specificities of the beauty and wellness segment}

Beauty, parapharmacy, and wellness retail present critical domain properties that distinguish them from standard consumer goods:

\paragraph{Suitability is an absolute safety constraint.} Unlike apparel or entertainment where an imperfect recommendation merely represents a minor aesthetic mismatch, cosmetic suitability is physiological. Formulations designed for oily or acne-prone skin containing high concentrations of comedogenic oils or aggressive exfoliants can cause acute inflammation, allergic flare-ups, or barrier damage when applied to dry or hyper-reactive skin. Consequently, negative contraindications must act as absolute disqualifications rather than soft score penalties.

\paragraph{Products operate as structured routines.} Skincare products are rarely applied in isolation; they function as a sequential, multi-step regimen (Cleansing $\rightarrow$ Toning $\rightarrow$ Targeted Treatment / Serum $\rightarrow$ Hydration / Moisturizing $\rightarrow$ Daytime Photoprotection). Recommending five cleansers simultaneously is functionally useless to a consumer; the system must assemble an ordered, balanced routine.

\paragraph{The semantic and vocabulary chasm.} Shoppers express concerns using symptomatic everyday vernacular (\emph{``My skin turns red after showering'', ``Dark patches appeared during pregnancy''}), whereas inventory records use biochemical dermatological terms (\emph{``erythema'', ``azelaic acid'', ``melasma'', ``sebum-regulating''}). Bridging this semantic divide is the core function of an intelligent adviser.

\paragraph{Sensitive cross-selling boundaries.} In wellness retail, topical cosmetics are frequently merchandised alongside ingestible food supplements. While complementary, cross-selling ingestible supplements must avoid unauthorized therapeutic or medical claims, requiring strict regulatory boundaries in agent dialogue.

\subsection{Opportunities and threats (SWOT analysis)}

Table~\ref{tab:swot} summarizes the strengths, weaknesses, opportunities, and threats associated with introducing an AI advisory recommender in the client's store.

\begin{table}[H]
\caption{SWOT analysis of the proposed solution}
\label{tab:swot}
\centering
\begin{tabularx}{\textwidth}{L{3.1cm} X}
\toprule
\textbf{Strengths} & Operates reliably on catalogues lacking structured tags; zero recurring inference cost for the deterministic scoring engine; strict explainability; zero-hallucination guarantee; modular architecture reusable across WooCommerce stores. \\
\midrule
\textbf{Weaknesses} & Initial rule extraction depends on description quality; conversational LLM tier subject to free-tier rate limitations unless scaled to paid infrastructure; initial absence of clinical dermatological validation. \\
\midrule
\textbf{Opportunities} & Reusable across all e-commerce clients of \companyname; algorithmic discovery of inventory gaps and supply shortages; substantial commercial differentiation against competitors relying on basic text search. \\
\midrule
\textbf{Threats} & Catalogue drift caused by newly added unlabelled inventory; sudden policy changes or deprecation of third-party LLM endpoints; evolving regulatory scrutiny around AI-generated health advice. \\
\bottomrule
\end{tabularx}
\end{table}

\section{Core concepts}

\subsection{Conversational agents and chatbots}

Conversational software has evolved from early rule-based pattern matchers (such as ELIZA or finite-state decision trees) to modern generative agents. In retail, standard chatbots follow fixed dialogue graphs that break when users formulate multi-faceted or unpredicted requests. An intelligent conversational agent, by contrast, dynamically interprets free-form language, tracks conversational state, queries internal tools, and synthesizes natural responses.

\subsection{Natural Language Processing (NLP)}

Natural Language Processing encompasses computational methods for analysing, parsing, and representing human language. Key operations relevant to e-commerce text mining include lexical normalisation (lowercasing, accent removal, punctuation stripping), stop-word filtering, n-gram tokenization, and named entity recognition. These techniques enable extracting structured semantic primitives from messy, marketing-heavy product copy.

\subsection{Large Language Models (LLMs)}

Large Language Models based on the Transformer architecture \cite{vaswani2017} and trained on massive web-scale corpora exhibit remarkable zero-shot and few-shot reasoning capabilities \cite{brown2020}. In this project, LLMs serve two distinct roles: converting ambiguous natural-language user queries into structured semantic profiles, and phrasing fluent, empathetic shopping advice grounded in verified product data.

\subsection{Information extraction from unstructured commercial text}

Deriving structured product attributes from unstructured catalog descriptions is an information extraction challenge. In e-commerce text mining, two primary paradigms exist:
\begin{itemize}
  \item \textbf{Rule-based pattern matching}: Leverages compiled regular expressions and domain lexicons. It is deterministic, fully inspectable, executes at sub-second speed at zero marginal cost, and is easily auditable. Because commercial cosmetic French copy is formulaic (consistently using constructions such as \emph{``destiné aux peaux mixtes à grasses'', ``lutte contre les taches brunes''}), pattern matching provides high precision.
  \item \textbf{Model-based extraction}: Uses generative language models to emit structured JSON schemas from free text. While adaptable to unexpected phrasings, it incurs per-call inference latency and API cost.
\end{itemize}

A key weighting refinement comes from information retrieval: **Inverse Document Frequency (IDF)** \cite{sparckjones1972, robertson2004}, which weights terms inversely to their catalogue frequency. Ubiquitous marketing words receive low weights, whereas rare identifying terms (such as specific botanical names) carry high discriminative power.

\subsection{Tool-augmented agents and Function Calling (ReAct)}

The ReAct framework \cite{yao2023react} and tool-calling protocols \cite{schick2023toolformer} allow language models to interface directly with external software tools, databases, and APIs. When presented with formal JSON schema signatures, the model emits structured function call requests rather than text. The execution engine executes the tool and supplies the resulting data back to the model context.

This introduces a crucial architectural boundary: the model does not attempt to remember or invent catalogue inventory; it acts as an intelligent orchestrator that translates customer needs into tool parameters, while deterministic scoring functions provide verified, real-time product data.

\subsection{Grounding and anti-hallucination guardrails}

Large language models are inherently vulnerable to hallucination \cite{ji2023survey}---generating factually inaccurate claims, non-existent product titles, or fictitious prices. In retail e-commerce, such hallucinations directly damage customer trust. 

Rather than relying purely on soft prompt instructions to curb hallucinations, the system implements **structural grounding**: customer-facing product recommendations are instantiated exclusively from validated database records returned by tool calls. If a model generates text referencing a product that was not returned by a tool, that product has no rendering path into the interface.

\subsection{Microservice containerisation with Docker}

Modern web architectures require containerisation to eliminate environment discrepancy across local development, testing, and production servers. Docker encapsulates the Python 3 runtime, FastAPI ASGI application, and database drivers into an isolated container orchestrated via Docker Compose.

\section{Study of existing solutions}

\subsection{Classical recommender algorithms}

The recommender systems literature \cite{ricci2015handbook, resnick1997} establishes three classical algorithmic families:
\begin{itemize}
  \item \textbf{Collaborative filtering}: Infers user preferences from past user-item interaction matrices \cite{sarwar2001}. It requires no semantic knowledge of product contents.
  \item \textbf{Content-based filtering}: Matches customer attribute profiles with structured product metadata \cite{lops2011}.
  \item \textbf{Hybrid recommenders}: Combine collaborative and content-based signals to balance novelty and accuracy \cite{burke2002}.
\end{itemize}

\subsection{The cold-start problem}

Collaborative filtering fails completely when user-item interaction matrices are empty or sparse---the notorious cold-start problem \cite{schein2002}. In the client's store, no customer purchase histories were linked to recommendation profiles, and new visitors arrive with zero prior interactions. This mathematically disqualifies collaborative filtering, making content-based filtering the only viable paradigm.

\subsection{Commercial WooCommerce tools and SaaS platforms}

Table~\ref{tab:comparison} contrasts existing market solutions against the client's operational requirements.

\begin{table}[H]
\caption{Comparison of existing approaches against project requirements}
\label{tab:comparison}
\centering
\begin{tabularx}{\textwidth}{L{3.0cm} X L{3.5cm}}
\toprule
\textbf{Category} & \textbf{Core mechanism} & \textbf{Fatal limitation on client store} \\
\midrule
WooCommerce recommendation plugins & Related products and cross-sells based on manual tags and categories & Assumes clean tags and taxonomies; the client store has zero tags and 6 populated skin types \\
\midrule
Product filter plugins & Faceted checkbox navigation over database attributes & Filters the same absent attributes; forces the shopper to diagnose their own skin \\
\midrule
Enterprise SaaS engines & Cloud behavioral matrix factorization & Requires massive traffic volume; high monthly subscription fees; unsuited for routine assembly \\
\midrule
Generic LLM chatbots & Text generation over ungrounded prompts & Prone to catastrophic hallucinations: inventing non-existent items, outdated prices, or out-of-stock SKUs \\
\bottomrule
\end{tabularx}
\end{table}

\section{Limitations of related work}

Three critical shortcomings disqualify existing commercial tools for the client:
\begin{enumerate}
  \item \textbf{Limited integration with internal relational data}: Generic AI chatbots have no real-time connection to WooCommerce inventory tables, leading to inaccurate stock reporting and mismatched pricing.
  \item \textbf{Prohibitive cost and reproducibility barriers}: Cloud enterprise SaaS platforms impose recurring subscription fees that small e-commerce merchants cannot sustain.
  \item \textbf{Hallucinations and lack of traceability}: An ungrounded generative model can confidently invent cosmetic formulations or make unsubstantiated medical claims, exposing the retailer to severe legal and reputational liability \cite{ji2023survey}.
\end{enumerate}

\section{Proposed solution}

To overcome these limitations, this project designs and implements a **Tool-Grounded Hybrid Advisory Recommender**. Figure~\ref{fig:rag_vs_tool_agent} contrasts standard conversational approaches with the proposed architecture.

\placeholderscreenshot{rag_vs_tool_agent.png}{Architectural comparison between ungrounded LLM chatbots and our Tool-Grounded ReAct Agent}{fig:rag_vs_tool_agent}

The core principles of the proposed solution are:
\begin{enumerate}
  \item \textbf{Automated information extraction before recommendation}: Rather than requiring hundreds of hours of manual catalog tagging, an automated pipeline parses unstructured descriptions into structured categorical attributes.
  \item \textbf{Deterministic selection, generative expression}: Candidate products and ordered routines are assembled exclusively by a deterministic mathematical scoring engine. The LLM is restricted to parsing customer intent and phrasing conversational responses; it has no programmatic path to invent products.
  \item \textbf{Tool-grounded ReAct loop}: The agent interacts with the inventory through three formal tools (\texttt{lister\_options}, \texttt{chercher\_produits}, and \texttt{construire\_routine}), guaranteeing that all recommended cards reflect real-time pricing and stock status.
\end{enumerate}

\section{Adopted methodology: CRISP-DM}

To guarantee engineering rigor and industrial reproducibility, the project was conducted according to the **CRISP-DM** (Cross-Industry Standard Process for Data Mining) methodology \cite{wirth2000crisp}. Figure~\ref{fig:crisp_dm} illustrates the cyclical stages of CRISP-DM adapted to this e-commerce project.

\placeholderscreenshot{crisp_dm_cycle.png}{The CRISP-DM methodology lifecycle applied to the product recommendation project}{fig:crisp_dm}

The six CRISP-DM phases were mapped as follows:
\begin{enumerate}
  \item \textbf{Business Understanding}: Formalising the client's business objective---transforming unstructured catalog browsing into an interactive cosmetic adviser that increases average order value and customer satisfaction.
  \item \textbf{Data Understanding}: Auditing the MariaDB database export, evaluating table relationships, assessing taxonomic sparsity, and quantifying textual completeness across 1{,}073 product records.
  \item \textbf{Data Preparation}: Cleansing HTML tags, normalising French text, engineering regex extraction rules, resolving commercial pack compositions, and generating structured JSON knowledge bases.
  \item \textbf{Modeling}: Designing the multi-attribute scoring algorithm, establishing routine assembly heuristics, implementing the FastAPI backend, and engineering the ReAct agent prompt and tool schema.
  \item \textbf{Evaluation}: Systematically validating the recommender against 22 representative customer profiles, measuring routine coherence, budget compliance, latency, and absence of hallucinations.
  \item \textbf{Deployment}: Containerising the service using Docker and Docker Compose, developing the custom WordPress plugin, and embedding shortcode interfaces into the live storefront.
\end{enumerate}

\section{Conclusion}

Personalised beauty advice requires content-based advisory recommendation rather than collaborative filtering. Because commercial off-the-shelf plugins fail when product attributes are missing, the automated extraction of structured data from raw text becomes the foundational entry point. By pairing deterministic scoring with an agentic ReAct loop under the CRISP-DM methodology, the system achieves the conversational fluency of generative AI with the strict reliability and safety of a deterministic database.

\clearpage
